% Part: normal-modal-logic
% Chapter: filtrations
% Section: more-filtrations

\documentclass[../../../include/open-logic-section]{subfiles}

\begin{document}

\olfileid{nml}{fil}{acc}

\olsection{ફિલ્ટ્રેશનો અને પ્રાપ્યતાના ગુણધર્મો}

જણાવ્યા મુજબ, સાર્વત્રિક, સિરિયલ અને સ્વવાચક નિદર્શનોનાં
ફિલ્ટ્રેશનો પણ હંમેશા અનુક્રમે સાર્વત્રિક, સિરિયલ અથવા સ્વવાચક
હોય છે. પરંતુ સંમિત અથવા પરંપરિત નિદર્શનનું દરેક ફિલ્ટ્રેશન
અનુક્રમે સંમિત અથવા પરંપરિત હોતું નથી. છતાં, કેટલાક કિસ્સામાં
આ ગુણધર્મ જળવાય તે રીતે ફિલ્ટ્રેશન વ્યાખ્યાયિત કરી શકાય છે.
એ માટે આપણે સૌથી સ્થૂલ ફિલ્ટ્રેશનની વ્યાખ્યા જેવી રીત લઈએ છીએ,
પણ~$R^*$ની વ્યાખ્યામાં વધારાની શરતો ઉમેરીએ છીએ. $\Gamma$
ઉપ-!!{formula}s હેઠળ બંધ માનો. નિદર્શન
$\mModel{M} =\tuple{W, R, V}$નાં જગતો $u$, $v$ વચ્ચેના
\olref{tab:Cn-filtrations}માંના સંબંધો~$C_i(u,v)$ જુઓ.
આ શરતોના સંયોજનો પરથી $R^*[u][v]$ વ્યાખ્યાયિત કરી શકાય છે.
દાખલા તરીકે, $R^*[u][v]$ જો અને માત્ર જો $C_1(u,v)$ હોય એવી
શરત રાખીએ, તો બરાબર સૌથી સ્થૂલ ફિલ્ટ્રેશન મળે છે. જો
$R^*[u][v]$ જો અને માત્ર જો $C_1(u,v)$ અને $C_2(u, v)$ બંને
હોય એમ રાખીએ, તો જુદું ફિલ્ટ્રેશન મળે છે. તે સૌથી સ્થૂલ
ફિલ્ટ્રેશન કરતાં ``વધુ સૂક્ષ્મ'' છે, કારણ કે એકલા $C_1(u,v)$
કરતાં $C_1(u,v)$ અને $C_2(u,v)$ બંનેને ઓછાં જગત-યુગ્મો
સંતોષે છે.

\begin{table}[ht]
  \centering
  \begin{tabular}{|ll|}
    \hline
    \multirow{2}{*}{$C_1(u,v)$:}
    \iftag{prvBox}{&
      જો $\Box!A \in \Gamma$ અને $\mSat{M}{\Box!A}[u]$ હોય,
      તો $\mSat{M}{!A}[v]$;\iftag{prvDiamond}{ અને}{}\\}{}
    \iftag{prvDiamond}{&
      જો $\Diamond!A \in \Gamma$ અને $\mSat{M}{!A}[v]$ હોય,
      તો $\mSat{M}{\Diamond!A}[u]$; \\}{}
    \hline
    \multirow{2}{*}{$C_2(u,v)$:}
    \iftag{prvBox}{&
      જો $\Box!A \in \Gamma$ અને $\mSat{M}{\Box!A}[v]$ હોય,
      તો $\mSat{M}{!A}[u]$;\iftag{prvDiamond}{ અને}{}\\}{}
    \iftag{prvDiamond}{&
      જો $\Diamond!A \in \Gamma$ અને $\mSat{M}{!A}[u]$ હોય,
      તો $\mSat{M}{\Diamond!A}[v]$; \\}{}
    \hline
    \multirow{2}{*}{$C_3(u,v)$:}
    \iftag{prvBox}{&
      જો $\Box!A \in \Gamma$ અને $\mSat{M}{\Box!A}[u]$ હોય,
      તો $\mSat{M}{\Box!A}[v]$;\iftag{prvDiamond}{ અને}{}\\}{}
    \iftag{prvDiamond}{&
      જો $\Diamond!A \in \Gamma$ અને $\mSat{M}{\Diamond!A}[v]$ હોય,
      તો $\mSat{M}{\Diamond!A}[u]$; \\}{}
    \hline
    \multirow{2}{*}{$C_4(u,v)$:}
    \iftag{prvBox}{&
      જો $\Box!A \in \Gamma$ અને $\mSat{M}{\Box!A}[v]$ હોય,
      તો $\mSat{M}{\Box!A}[u]$;\iftag{prvDiamond}{ અને}{}\\}{}
    \iftag{prvDiamond}{&
      જો $\Diamond!A \in \Gamma$ અને $\mSat{M}{\Diamond!A}[u]$ હોય,
      તો $\mSat{M}{\Diamond!A}[v]$; \\}{}
    \hline
    \end{tabular}
  \caption{ફિલ્ટ્રેશનો વ્યાખ્યાયિત કરવા માટે જગતો પરની શરતો.}
  \ollabel{tab:Cn-filtrations}
\end{table}

\begin{thm}\ollabel{thm:more-filtrations}
  $\mModel{M} =\tuple{W,R,V}$ નિદર્શન હોય, $\Gamma$ ઉપ-!!{formula}s
  હેઠળ બંધ હોય અને $W^*$ તથા $V^*$ની વ્યાખ્યા
  \olref[fil]{defn:filtration} મુજબ હોય. ત્યારે:
  \begin{enumerate}
  \item ધારો કે $R^*[u][v]$ જો અને માત્ર જો $C_1(u, v) \land
    C_2(u,v)$. ત્યારે $R^*$ સંમિત છે; અને $\mModel{M}$ સંમિત હોય
    તો $\mModel{M^*} = \tuple{W^*,R^*,V^*}$ ફિલ્ટ્રેશન છે.
  \item ધારો કે $R^*[u][v]$ જો અને માત્ર જો $C_1(u, v)
    \land C_3(u,v)$. ત્યારે $R^*$ પરંપરિત છે; અને
    $\mModel{M}$ પરંપરિત હોય તો
    $\mModel{M^*}=\tuple{W^*,R^*,V^*}$ ફિલ્ટ્રેશન છે.
  \item ધારો કે $R^*[u][v]$ જો અને માત્ર જો $C_1(u, v) \land C_2(u,v)
    \land C_3(u,v) \land C_4(u,v)$. ત્યારે $R^*$ સંમિત અને
    પરંપરિત છે; અને $\mModel{M}$ સંમિત તથા પરંપરિત હોય તો
    $\mModel{M^*}=\tuple{W^*,R^*,V^*}$ ફિલ્ટ્રેશન છે.
  \item ધારો કે $R^*$ એ રીતે વ્યાખ્યાયિત છે કે $R^*[u][v]$ જો
    અને માત્ર જો $C_1(u,
    v) \land C_3(u,v) \land C_4(u,v)$. ત્યારે $R^*$ પરંપરિત
    અને યુક્લિડિયન છે; અને $\mModel{M}$ પરંપરિત તથા યુક્લિડિયન
    હોય તો $\mModel{M^*}=\tuple{W^*,R^*,V^*}$ ફિલ્ટ્રેશન છે.
  \end{enumerate}
\end{thm}

\begin{proof}
  \begin{enumerate}
    \item $R^*$ સંમિત છે તે તરત જ દેખાય છે, કારણ કે $C_1(u,v)
      \Leftrightarrow C_2(v,u)$ અને $C_2(u,v) \Leftrightarrow
      C_1(v,u)$. તેથી $\mModel{M}$ સંમિત હોય તો
      $\mModel{M^*}$ એ~$\Gamma$ દ્વારા ફિલ્ટ્રેશન છે તે
      બતાવવાનું રહે છે. $C_1(u,v)$ શરતથી
      \iftag{prvBox}{%
        \iftag{prvDiamond}{%
          \olref[fil]{defn:filtration-R2} અને
          \olref[fil]{defn:filtration-R3}, એટલે
          \olref[fil]{defn:filtration}ની શરતો}{%
          \olref[fil]{defn:filtration}%
          \olref[fil]{defn:filtration-R2} શરત}}{%
        \olref[fil]{defn:filtration}%
        \olref[fil]{defn:filtration-R3} શરત}
      પૂરી થાય છે. આમ માત્ર
      \olref[fil]{defn:filtration}%
      \olref[fil]{defn:filtration-R1} ચકાસવી રહે છે, એટલે કે
      $Ruv$ હોય તો $R^*[u][v]$.

      $Ruv$ ધારો. $R^*[u][v]$ બતાવવા $C_1(u,v)$ અને
      $C_2(u,v)$ સ્થાપિત કરવાં પડે. $C_1$ માટે:
      \iftag{prvBox}{જો $\Box!A
        \in\Gamma$ અને $\mSat{M}{\Box!A}[u]$ હોય, તો
        $\mSat{M}{!A}[v]$ પણ (કારણ કે $Ruv$).
        \iftag{prvDiamond}{તેવી જ રીતે, }{}}{}%
      \iftag{prvDiamond}{જો $\Diamond!A \in \Gamma$ અને
        $\mSat{M}{!A}[v]$ હોય, તો $Ruv$ હોવાથી
        $\mSat{M}{\Diamond!A}[u]$.}{}
      $C_2$ માટે:
      \iftag{prvBox}{જો $\Box!A \in \Gamma$ અને
        $\mSat{M}{\Box!A}[v]$ હોય, તો સંમિતતાથી $Ruv$માંથી
        $Rvu$ મળે છે; એટલે $\mSat{M}{!A}[u]$.
        \iftag{prvDiamond}{તેવી જ રીતે, }{}}{}%
      \iftag{prvDiamond}{જો $\Diamond!A \in\Gamma$ અને
        $\mSat{M}{!A}[u]$ હોય, તો $\mSat{M}{\Diamond!A}[v]$
        (કારણ કે સંમિતતાથી $Rvu$).}{}
    \item અભ્યાસ માટે છોડ્યું છે.
    \item અભ્યાસ માટે છોડ્યું છે.
    \item અભ્યાસ માટે છોડ્યું છે.
  \end{enumerate}
\end{proof}

\begin{prob}
  \olref[nml][fil][acc]{thm:more-filtrations}ની સાબિતી પૂરી કરો.
\end{prob}

\end{document}
